
In this section, we present the evaluation dataset and baselines and afterwards discuss our main results.
\subsection{Datasets}
\label{sec:data}
Evaluation data were obtained from the open-access PK-DB repository, which aggregates pharmacokinetic measurements from clinical and pre-clinical studies~\citep{grzegorzewski2021pkdb}.
%
From PK-DB we curated plasma concentration--time measurements from a Phase~I study in healthy volunteers investigating the safety and dosing of a combination of drugs and their metabolites~\citep{lenuzza2016safety}.
%
In total, the dataset comprises 18 compounds (parent drugs and primary metabolites) with 2019 valid concentration--time observations.
%
Sampling frequencies vary substantially both within and across compounds, resulting in irregularly sampled time series.
%
Each compound is measured in 10 individuals, with up to 15 observations per individual. 
We also included two additional datasets on indometacin and theophylline pharmacokinetics (\texttt{Indometh} and \texttt{Theoph} from the R \texttt{datasets} package) \citep{Rsoftware}.

\subsection{Baselines and Ablations}
\label{sec:baselines}

We compare our proposed Prior-Fitted Functional Flows (PFF) against a classical pharmacometric baseline, nonlinear mixed-effects modeling (NLME), as well as two neural baselines: NODE-PK and the AICMET Transformer. The NLME baseline is fitted study-by-study for each compound using standard estimation procedures (one- or two-compartment models with first-order absorption, selected by AIC), including inter-individual variability and combined error models. Metabolites were excluded as they lack explicit dosing inputs. In contrast, all neural models are pre-trained on our synthetic datasets and evaluated in-context on empirical studies without fine-tuning. NODE-PK extends latent neural ODEs~\citep{rubanova2019latent} with explicit dosing handling via an RNN encoder. But note that it is not capable of study-level aggregation~\citep{lu2021deep}. The AICMET Transformer~\citep{marin2025amortized} serves as a hierarchical neural process baseline, capturing both population-level and individual-level effects via separate encoding variables. 
\subsection{Metrics}

\paragraph{Forecasting evaluation} The performance of the past-conditioned models were assessed using a leave-one-individual-out setting: predictions for subject $i$ were conditioned on the study context $\mathcal{S}$, dosing $\mathbf{u}^i$, and the first four PK samples $\mathcal{D}^i_{\text{C}}$, then evaluated on the remaining samples $\mathcal{D}^i_{\text{T}}$ using log-RMSE.\\ 


\paragraph{Generative evaluation} We evaluate time series sample quality in two separate settings: an \emph{empirical} setting and a \emph{synthetic} setting. 
%
In both cases, the goal is to compare a reference set of trajectories against a corresponding set of trajectories generated by the model under the same study-level conditioning, but the way these two sets are constructed differs across the two settings.
%
For the empirical evaluation, we apply a leave-one-out protocol across all empirical studies and all available drugs. 
%
In each study, one individual is removed from the study context, the model is conditioned on the remaining individuals, and one trajectory is sampled for the held-out individual. 
%
Repeating this procedure over all individuals in all studies yields a generated empirical set whose cardinality matches that of the full empirical dataset, since each empirical individual is held out and resampled exactly once. 
%
The resulting model-generated trajectories are then pooled across studies and compared against the pooled empirical trajectories from the original dataset.
%
For the synthetic evaluation, no leave-one-out construction is needed. 
%
Instead, each synthetic study provides a study context together with a reference population of target individuals sampled directly from the simulator. 
%
Conditioned on the same context, the model generates a matching synthetic population, and these model-generated trajectories are compared directly against the corresponding oracle trajectories.
%
After constructing these empirical and synthetic comparison pairs, we quantify discrepancy in both settings using two complementary two-sample metrics: 
%
(i) the area under the ROC curve (AUC) of a classifier trained to distinguish reference from generated trajectories, where values closer to $0.5$ indicate better sample fidelity, and 
%
(ii) a maximum mean discrepancy (MMD) computed with the signature kernel of \citet{kiraly2019kernels}, which measures distributional discrepancy between the two sets of irregular trajectories in path space.


Furthermore, the generative capability of the models were assessed via visual predictive checks (VPC), comparing simulation percentiles against observed concentrations~\citep{holford2005}. Full implementation details are provided in the Supplementary Material.


\begin{table}[t]
\centering
\caption{\textbf{Sample Fidelity Scores} Lower is better for MMD, while AUC values closer to $0.5$ indicate better agreement.}
\label{tab:sample_fidelity}
\scriptsize
\begin{tabular}{lcc}
\toprule
\textbf{Model} & \textbf{AUC} & \textbf{MMD} \\
\midrule
\multicolumn{3}{c}{\textbf{Empirical}} \\
\cmidrule(lr){1-3}
AICMET  & 	
0.577 & 	
-0.001 \\
FlowPK & 	
\textbf{0.567} & 	
\textbf{-0.003} \\
\addlinespace
\multicolumn{3}{c}{\textbf{Synthetic}} \\
\cmidrule(lr){1-3}
AICMET  & 0.541 & 	
0.002 \\
FlowPK & \textbf{0.508} & 	
\textbf{0.001} \\
\bottomrule
\end{tabular}
\end{table}

\subsection{Discussion of numerical results}

Table~\ref{tab:main_results} summarizes the quantitative forecasting performance. \texttt{PFF} achieves the lowest or on par log-RMSE across the evaluated compounds in the majority of cases. This predictive accuracy and robust uncertainty calibration are illustrated in Figure~\ref{fig:prediction_row}, which illustrates inferred future trajectories for three partially observed individuals. Furthermore, Figure~\ref{fig:vpc_only} confirms the model's generative capacity for zero-shot population synthesis; the Visual Predictive Checks demonstrate that \texttt{PFF}'s simulated percentiles tightly align with the empirical data distributions of entire study cohorts.
In Table~\ref{tab:sample_fidelity} we present the results for generated sample fidelity compared against the other generative model for Pharmacokinetic AICMET. Our model outperforms in both metrics the previous competing generative model.
%We compare a \emph{classical} pharmacometric baseline against two neural baselines—NODE-PK and the AICMET Transformer—and our proposed \emph{Prior Fitted Functional Flows} (PFFF). The classical baseline is nonlinear mixed-effects modelling (NLME), which we fit \emph{study by study}, i.e.\ separately for each compound-specific dataset using standard estimation procedures. In contrast, all neural models—NODE-PK, AICMET-Transformer, and PFFF—are \emph{pre-trained on our simulated PK datasets} and subsequently \emph{evaluated in context on the empirical studies}, without any additional per-study parameter updates or fine-tuning.

%\paragraph{Nonlinear Mixed-Effect (NLME) Modelling.} For each parent compound, we fitted one- and two-compartment NLME models (with first-order oral absorption for orally dosed compounds), including inter-individual variability on all parameters and a combined additive/proportional error model. Model order was selected by AIC and used as the compound-specific baseline. Metabolites were not modelled explicitly, since their effective input depends on the parent PK and cannot be treated as an observed dosing process within this setup.

%\paragraph{Neural ODEs (NODE-PK).} The NODE-PK model follows \citet{lu2021deep} and extends the latent neural ODE of \citet{rubanova2019latent} with explicit handling of dosing. An RNN encoder produces a latent initial condition for each individual, which is then evolved through neural ODE dynamics between observation times. NODE-PK does not perform study-level aggregation: individuals are treated independently, and predictions rely solely on their own past observations.

%\paragraph{AICMET Transformer.} As a neural baseline, we compare against AICMET \cite{marin2025amortized}, an amortized hierarchical neural process in which two encoding variables capture population-level (fixed) and individual-level (random) effects.

%\paragraph{Evaluation Metrics.} For each compound, predictive performance was assessed in a leave-one-individual-out setting: predictions for subject $i$ were conditioned on the study context $\mathcal{S}$, dosing $\mathbf{u}^i$, and the first four PK samples $\mathcal{D}^i_{\text{C}}$, and evaluated on the remaining samples $\mathcal{D}^i_{\text{T}}$ using log-RMSE. Generative capability was assessed via a \emph{visual predictive check} (VPC), comparing simulation percentiles against observed concentrations \citep{holford2005}. Our models 

%To evaluate generative capability, we performed a simulation-based graphical diagnostic known as the \emph{visual predictive check} (VPC), a standard tool in pharmacometrics \citep{holford2005}. This diagnostic tool compares the variability of simulated plasma concentration profiles within a study via an overlay of simulation percentiles and observed data.



%\paragraph{Implementation Details.} The Supplementary Material contains the full implementation details.


